% GENERATED FILE. Edit canonical lessons, then run npm run generate:books.
% canonical-source-hash: fnv1a64:2a35205b425b5437
% canonical-lessons: TE-C340-taksanam, TE-C340-anukokunda, TE-C340-unnattundi, TE-C340-tira, TE-C340-melkonu

\chapter{Out of the Blue}
\label{ch:te-out-of-the-blue}
\hlchaptermodality{340}

\begin{chapteropening}
\textbf{By the end of this chapter:} \emph{I can say instantly, unexpectedly, out of the blue, in the end and to wake up.}

Complete the last lesson of chapter 340: I can say instantly, unexpectedly, out of the blue, in the end and to wake up.
\end{chapteropening}

\section[\texorpdfstring{\te{తక్షణం} (takṣaṇaṁ)}{తక్షణం (takṣaṇaṁ)}]{\te{తక్షణం} (takṣaṇaṁ) — instantly; at once}

\label{lesson:TE-C340-taksanam}



\begin{quote}
Before the new one: say the Telugu for a march, then the Telugu for to march against.
\end{quote}

\subsection*{You'll want to know: \te{తక్షణం}}
\textbf{\te{తక్షణం}} — \emph{takṣaṇaṁ} — \textquotedblleft{}instantly; at once\textquotedblright{}.

\begin{quote}
\te{సాక్షి} \te{తక్షణం} \te{పోలీసులకు} \te{ఫోను} \te{చేశాడు}.
\end{quote}

\begin{quote}
\emph{sākṣi takṣaṇaṁ pōlīsulaku phōnu cēśāḍu.}
\end{quote}

\begin{quote}
The witness phoned the police at once.
\end{quote}

Literally that very moment: \te{తత్} and \te{క్షణం}. \te{తక్షణమే} is stronger still. Where \te{ఈలోగా}, meanwhile, leaves a little time, \te{తక్షణం} leaves none.

\begin{cousinweb}
\textbf{\te{తక్షణం}} (\emph{takṣaṇaṁ}) is borrowed from Sanskrit \emph{tatkṣaṇam}, \textquotedblleft{}at that very moment\textquotedblright{}, from \emph{tad}, \textquotedblleft{}that\textquotedblright{}, and \emph{kṣaṇa}, \textquotedblleft{}a moment\textquotedblright{} (CDIAL 5640).
\end{cousinweb}

\begin{grammarlens}[title={What you've built}]
One more word for telling what happened.
\end{grammarlens}

\subsection*{Your turn}
\begin{itemize}
\raggedright
  \item \emph{Say it:} \emph{takṣaṇaṁ}
  \item \emph{Say it:} \emph{takṣaṇaṁ}, once more
  \item \emph{Say it:} \emph{sākṣi takṣaṇaṁ pōlīsulaku phōnu cēśāḍu.} — The witness phoned the police at once.
  \item \emph{Recall:} say the Telugu for a march, then the Telugu for to march against, then say \emph{takṣaṇaṁ} again
\end{itemize}

\begin{tcolorbox}[breakable,colback=teal!4,colframe=teal!35!black,title={Before you move on}]
What does \te{తక్షణం} mean? (\textquotedblleft{}Instantly; at once\textquotedblright{}.) Say it once more.
\end{tcolorbox}
\section[\texorpdfstring{\te{అనుకోకుండా} (anukōkuṇḍā)}{అనుకోకుండా (anukōkuṇḍā)}]{\te{అనుకోకుండా} (anukōkuṇḍā) — unexpectedly}

\label{lesson:TE-C340-anukokunda}



\begin{quote}
Before the new one: say the Telugu for to march against, then the Telugu for instantly.
\end{quote}

\subsection*{You'll want to know: \te{అనుకోకుండా}}
\textbf{\te{అనుకోకుండా}} — \emph{anukōkuṇḍā} — \textquotedblleft{}unexpectedly\textquotedblright{}.

\begin{quote}
\te{బజారులో} \te{అనుకోకుండా} \te{మా} \te{మాస్టారు} \te{కనిపించారు}.
\end{quote}

\begin{quote}
\emph{bajārulō anukōkuṇḍā mā māsṭāru kanipiñcāru.}
\end{quote}

\begin{quote}
I unexpectedly ran into our teacher in the market.
\end{quote}

The negative of \te{అనుకో}, to think: without thinking of it.

\begin{cousinweb}
\textbf{\te{అనుకోకుండా}} (\emph{anukōkuṇḍā}) is native Telugu: \emph{anukonu}, \textquotedblleft{}to think\textquotedblright{} (DEDR 868), with \emph{-kuṇḍā}, \textquotedblleft{}without -ing\textquotedblright{}: without having thought of it.
\end{cousinweb}

\begin{grammarlens}[title={What you've built}]
One more word for telling what happened.
\end{grammarlens}

\subsection*{Your turn}
\begin{itemize}
\raggedright
  \item \emph{Say it:} \emph{anukōkuṇḍā}
  \item \emph{Say it:} \emph{anukōkuṇḍā}, once more
  \item \emph{Say it:} \emph{bajārulō anukōkuṇḍā mā māsṭāru kanipiñcāru.} — I unexpectedly ran into our teacher in the market.
  \item \emph{Recall:} say the Telugu for to march against, then the Telugu for instantly, then say \emph{anukōkuṇḍā} again
\end{itemize}

\begin{tcolorbox}[breakable,colback=teal!4,colframe=teal!35!black,title={Before you move on}]
What does \te{అనుకోకుండా} mean? (\textquotedblleft{}Unexpectedly\textquotedblright{}.) Say it once more.
\end{tcolorbox}
\section[\texorpdfstring{\te{ఉన్నట్టుండి} (unnaṭṭuṇḍi)}{ఉన్నట్టుండి (unnaṭṭuṇḍi)}]{\te{ఉన్నట్టుండి} (unnaṭṭuṇḍi) — out of the blue}

\label{lesson:TE-C340-unnattundi}



\begin{quote}
Before the new one: say the Telugu for instantly, then the Telugu for unexpectedly.
\end{quote}

\subsection*{You'll want to know: \te{ఉన్నట్టుండి}}
\textbf{\te{ఉన్నట్టుండి}} — \emph{unnaṭṭuṇḍi} — \textquotedblleft{}out of the blue\textquotedblright{}.

\begin{quote}
\te{ఉన్నట్టుండి} \te{వర్షం} \te{మొదలైంది}.
\end{quote}

\begin{quote}
\emph{unnaṭṭuṇḍi varṣaṁ modalaindi.}
\end{quote}

\begin{quote}
Out of the blue it started to rain.
\end{quote}

Literally while being as it was: all was calm, and then it was not.

\begin{cousinweb}
\textbf{\te{ఉన్నట్టుండి}} (\emph{unnaṭṭuṇḍi}) is native Telugu, built twice on \emph{uṇḍu}, \textquotedblleft{}to be, exist\textquotedblright{} (DEDR 697): while things were as they were.
\end{cousinweb}

\begin{grammarlens}[title={What you've built}]
One more word for telling what happened.
\end{grammarlens}

\subsection*{Your turn}
\begin{itemize}
\raggedright
  \item \emph{Say it:} \emph{unnaṭṭuṇḍi}
  \item \emph{Say it:} \emph{unnaṭṭuṇḍi}, once more
  \item \emph{Say it:} \emph{unnaṭṭuṇḍi varṣaṁ modalaindi.} — Out of the blue it started to rain.
  \item \emph{Recall:} say the Telugu for instantly, then the Telugu for unexpectedly, then say \emph{unnaṭṭuṇḍi} again
\end{itemize}

\begin{tcolorbox}[breakable,colback=teal!4,colframe=teal!35!black,title={Before you move on}]
What does \te{ఉన్నట్టుండి} mean? (\textquotedblleft{}Out of the blue\textquotedblright{}.) Say it once more.
\end{tcolorbox}
\section[\texorpdfstring{\te{తీరా} (tīrā)}{తీరా (tīrā)}]{\te{తీరా} (tīrā) — in the end; when it came to it}

\label{lesson:TE-C340-tira}



\begin{quote}
Before the new one: say the Telugu for unexpectedly, then the Telugu for out of the blue.
\end{quote}

\subsection*{You'll want to know: \te{తీరా}}
\textbf{\te{తీరా}} — \emph{tīrā} — \textquotedblleft{}in the end; when it came to it\textquotedblright{}.

\begin{quote}
\te{చాలా} \te{రోజులు} \te{వెతికాము}, \te{తీరా} \te{అది} \te{ఇంట్లోనే} \te{ఉంది}.
\end{quote}

\begin{quote}
\emph{cālā rōjulu vetikāmu, tīrā adi iṇṭlōnē undi.}
\end{quote}

\begin{quote}
We searched for days, and in the end it was right at home.
\end{quote}

A spoken word that brings a surprise ending, often an ironic one.

\begin{cousinweb}
\textbf{\te{తీరా}} (\emph{tīrā}) is native Dravidian, built on \emph{tīru}, \textquotedblleft{}to be finished, concluded\textquotedblright{} (DEDR 3278): when everything was done.
\end{cousinweb}

\begin{grammarlens}[title={What you've built}]
One more word for telling what happened.
\end{grammarlens}

\subsection*{Your turn}
\begin{itemize}
\raggedright
  \item \emph{Say it:} \emph{tīrā}
  \item \emph{Say it:} \emph{tīrā}, once more
  \item \emph{Say it:} \emph{cālā rōjulu vetikāmu, tīrā adi iṇṭlōnē undi.} — We searched for days, and in the end it was right at home.
  \item \emph{Recall:} say the Telugu for unexpectedly, then the Telugu for out of the blue, then say \emph{tīrā} again
\end{itemize}

\begin{tcolorbox}[breakable,colback=teal!4,colframe=teal!35!black,title={Before you move on}]
What does \te{తీరా} mean? (\textquotedblleft{}In the end; when it came to it\textquotedblright{}.) Say it once more.
\end{tcolorbox}
\section[\texorpdfstring{\te{మేల్కొను} (mēlkonu)}{మేల్కొను (mēlkonu)}]{\te{మేల్కొను} (mēlkonu) — to wake up; to awaken}

\label{lesson:TE-C340-melkonu}



\begin{quote}
Before the new one: say the Telugu for out of the blue, then the Telugu for in the end.
\end{quote}

\subsection*{You'll want to know: \te{మేల్కొను}}
\textbf{\te{మేల్కొను}} — \emph{mēlkonu} — \textquotedblleft{}to wake up; to awaken\textquotedblright{}.

\begin{quote}
\te{నేను} \te{పొద్దున్నే} \te{ఆరు} \te{గంటలకు} \te{మేల్కొన్నాను}.
\end{quote}

\begin{quote}
\emph{nēnu poddunnē āru gaṇṭalaku mēlkonnānu.}
\end{quote}

\begin{quote}
I woke up at six in the morning.
\end{quote}

Past (I …): \te{మేల్కొన్నాను}, \emph{mēlkonnānu}.

A slightly bookish partner of \te{నిద్ర} \te{లేచు}.

\begin{cousinweb}
\textbf{\te{మేల్కొను}} (\emph{mēlkonu}) is native Dravidian: DEDR 5086 lists Telugu \emph{mēlu-konu}, \textquotedblleft{}to wake up, get up\textquotedblright{}, under the word for \textquotedblleft{}above, high\textquotedblright{}.
\end{cousinweb}

\begin{grammarlens}[title={What you've built}]
One more word for telling what happened.
\end{grammarlens}

\subsection*{Your turn}
\begin{itemize}
\raggedright
  \item \emph{Say it:} \emph{mēlkonu}
  \item \emph{Say it:} \emph{mēlkonu}, once more
  \item \emph{Say it:} \emph{nēnu poddunnē āru gaṇṭalaku mēlkonnānu.} — I woke up at six in the morning.
  \item \emph{Recall:} say the Telugu for out of the blue, then the Telugu for in the end, then say \emph{mēlkonu} again
\end{itemize}

\begin{tcolorbox}[breakable,colback=teal!4,colframe=teal!35!black,title={Before you move on}]
What does \te{మేల్కొను} mean? (\textquotedblleft{}To wake up; to awaken\textquotedblright{}.) Say it once more.
\end{tcolorbox}
